\section{Methodology}
\label{sec:method}

We design a controlled experiment to isolate the effect of CoT on output convergence. The key idea is simple: prompt the same models on the same questions under two conditions (with and without CoT), generate multiple samples per condition, and compare the degree of agreement.

\subsection{Experimental Design}

\para{Prompting conditions.} We compare two conditions:
\begin{itemize}[leftmargin=*,itemsep=0pt,topsep=0pt]
    \item \textbf{Direct}: The model is instructed to provide only the final answer, with no reasoning.
    \item \textbf{Zero-shot CoT}: The prompt includes ``Let's think step by step'' and requests reasoning before the final answer, following \citet{kojima2022large}.
\end{itemize}
We use zero-shot CoT rather than few-shot to avoid exemplar-specific confounds. Max token limits are set to 150 (direct) and 600 (CoT) to accommodate reasoning chains.

\para{Models.} We use three models from the GPT-4.1 family, accessed via the OpenAI API:
\begin{itemize}[leftmargin=*,itemsep=0pt,topsep=0pt]
    \item \gptlarge --- the largest model in the family
    \item \gptmini --- a medium-sized variant
    \item \gptnano --- the smallest variant
\end{itemize}
Using models from the same family at different scales controls for training data and architecture while varying capability, allowing us to study how model size modulates convergence effects.

\para{Sampling.} For each question $\times$ model $\times$ condition, we generate $K = 5$ independent samples at temperature $T = 0.7$. Temperature $> 0$ is necessary to observe output diversity (deterministic decoding trivially yields perfect agreement). We follow \citet{wang2022self} in using $T = 0.7$, which provides a balance between diversity and coherence.

\subsection{Datasets}

We sample 30 questions from each of three benchmarks, selected to span different reasoning domains and answer formats:

\begin{table}[h]
\centering
\begin{tabular}{@{}llll@{}}
\toprule
\textbf{Dataset} & \textbf{Task Type} & \textbf{Answer Format} & \textbf{Split} \\
\midrule
\gsm \citep{cobbe2021gsm8k} & Math reasoning & Numeric (open) & Test \\
\csqa \citep{talmor2019commonsenseqa} & Commonsense & 5-choice (A--E) & Validation \\
\arc \citep{clark2018arc} & Science & 4-choice (A--D) & Test \\
\bottomrule
\end{tabular}
\caption{Datasets used in our experiments. We sample 30 questions per dataset (seed = 42), yielding 90 questions total.}
\label{tab:datasets}
\end{table}

The distinction between \gsm (open-ended numeric answers) and the multiple-choice datasets (\csqa, \arc) turns out to be critical for understanding CoT's convergence effects, as we show in \secref{sec:results}.

\subsection{Metrics}
\label{sec:metrics}

\para{Answer agreement rate (\aar).} For a given question, model, and condition, we compute the pairwise agreement among the $K = 5$ samples. Formally, for samples $\{a_1, \ldots, a_K\}$:
\begin{equation}
\text{AAR} = \frac{1}{\binom{K}{2}} \sum_{i < j} \mathbf{1}[a_i = a_j]
\label{eq:aar}
\end{equation}
An \aar of 1.0 means all 5 samples produced the same answer; lower values indicate more diverse outputs.

\para{Cross-model pairwise agreement.} For a given question and condition, we compute pairwise agreement across all $K \times K = 25$ sample pairs between two models. This measures whether different models converge on the same answers.

\para{Output entropy.} We compute the Shannon entropy of the answer distribution per question:
\begin{equation}
H = -\sum_{a \in \gA} p(a) \log_2 p(a)
\label{eq:entropy}
\end{equation}
where $\gA$ is the set of distinct answers and $p(a)$ is the empirical frequency. Lower entropy indicates more convergent outputs.

\para{Accuracy.} We report standard accuracy (fraction of samples matching the gold answer) to contextualize convergence with correctness.

\subsection{Statistical Analysis}

We use the Wilcoxon signed-rank test to compare CoT versus direct conditions, paired by question. This non-parametric test is appropriate because our metrics (\aar, agreement rates) are bounded and potentially non-normal. We report Cohen's $d$ as the effect size, computed on the paired differences. All tests use significance level $\alpha = 0.05$. We conduct 2,700 total API calls: 3 models $\times$ 90 questions $\times$ 2 conditions $\times$ 5 samples.
